\chapter{Discussion and Future Work}
\label{chap:conclusion}
\label{sec:discussion}

We proposed DART, a denoising-score-assisted Rao detector that combines
classical weak-target detection with score-based background learning, and
its spatially adapted extension DARTS. The learned background model is
reusable across target signatures, and in the affine Gaussian case DART
reduces to a diagonally loaded adaptive matched filter, explaining its
finite-sample robustness.

Many directions remain open. First, stronger score-learning objectives may
improve detection; in particular, flow-matching or multi-noise
diffusion-style alternatives could learn more accurate scores. Second,
DARTS contains several manually selected hyperparameters, including the
network size, neighborhood size, number of neighbors, and
the shrinkage parameter $\lambda$; automatic tuning using only secondary
data is an important practical next step. Finally, robustness to target
leakage in the secondary data and a systematic ablation of the whitening
front-end remain to be studied.

Overall, DART shows that score-based learning can serve as a focused
replacement for the unknown background model inside a classical detector,
rather than as an end-to-end learned classifier, extending adaptive
detection beyond restrictive Gaussian assumptions while preserving
interpretability and calibration.
